% Generated by tools/tpt document-types refresh from src/common/document-types/server-sung-guide-v1/main.template.tex
% and this leaf's publication.toml [document] table. Do not edit.
\documentclass[11pt]{article}
\usepackage[margin=0.75in]{geometry}
\usepackage[T1]{fontenc}
\usepackage[utf8]{inputenc}
\usepackage{lmodern}
\usepackage{microtype}
\usepackage{array}
\usepackage{booktabs}
\usepackage{longtable}
\usepackage{tabularx}
\usepackage{enumitem}
\usepackage{needspace}
\usepackage{multicol}
\usepackage{xcolor}
\usepackage{hyperref}
\usepackage[most]{tcolorbox}
\usepackage{tikz}
\usetikzlibrary{arrows.meta,calc,positioning,shapes.geometric}

% Semantic color names are retained for compatibility with the weekly files,
% but every value is pure black or white for monochrome printing.
\definecolor{rubricred}{gray}{0}
\definecolor{linkblue}{gray}{0}
\definecolor{softblue}{gray}{1}
\definecolor{softgold}{gray}{1}
\definecolor{softgray}{gray}{1}

\hypersetup{
  hidelinks,
  pdfborder={0 0 0}
}

\setlength{\parindent}{0pt}
\setlength{\parskip}{0.55em}
\setlength{\tabcolsep}{3pt}
\setlist[itemize]{leftmargin=1.25em, itemsep=0.2em}
\renewcommand{\arraystretch}{1.22}

\newcommand{\weektitle}[3]{
  \begin{center}
    {\Large\bfseries #1}\\[0.2em]
    {\large #2}\\[0.2em]
    {\small #3}
  \end{center}
  \vspace{0.5em}
}

\newcommand{\properrow}[4]{\textbf{#1} & #2 & #3 & #4\\}
\newcommand{\properrefs}[1]{\textnormal{(\textit{#1})}}
\newcommand{\sourceurl}[2]{\href{#1}{#2}}
\newcommand{\latin}[1]{\textit{#1}}
\newcommand{\rubric}[1]{\textbf{#1}}

\newenvironment{properstable}
{\footnotesize\begin{longtable}{>{\raggedright\arraybackslash}p{0.15\linewidth}>{\raggedright\arraybackslash}p{0.25\linewidth}>{\raggedright\arraybackslash}p{0.18\linewidth}>{\raggedright\arraybackslash}p{0.34\linewidth}}
\textbf{Proper} & \textbf{TLM text / reference} & \textbf{Scriptural axis} & \textbf{Connection}\\
\midrule
\endhead}
{\bottomrule\end{longtable}}

\newenvironment{sensestable}
{\vspace{-0.35em}\begingroup\footnotesize\renewcommand{\arraystretch}{1.10}%
\noindent\hspace*{\dimexpr0.04\linewidth-4\tabcolsep\relax}%
\begin{tabular}{>{\bfseries\raggedright\arraybackslash}p{0.15\linewidth}>{\raggedright\arraybackslash}p{\dimexpr0.77\linewidth+4\tabcolsep\relax}}
\textbf{Sense} & \textbf{Synthesis}\\
\midrule}
{\bottomrule\end{tabular}\par\endgroup}

\newenvironment{connectionstable}
{\small\begin{longtable}{>{\bfseries\raggedright\arraybackslash}p{0.18\linewidth}>{\raggedright\arraybackslash}p{0.74\linewidth}}
\toprule
\textbf{Thread} & \textbf{Synthesis}\\
\midrule
\endhead}
{\bottomrule\end{longtable}}

\newenvironment{witnessstable}
{\small\begin{longtable}{>{\raggedright\arraybackslash}p{0.2\linewidth}>{\raggedright\arraybackslash}p{0.34\linewidth}>{\raggedright\arraybackslash}p{0.38\linewidth}}
\toprule
\textbf{Witness} & \textbf{Text / source status} & \textbf{Interpretive contribution}\\
\midrule
\endhead}
{\bottomrule\end{longtable}}

\newcommand{\witnessrow}[3]{#1 & #2 & #3\\}
\newcommand{\sectionguard}{\Needspace{7\baselineskip}}

% Keep the opening survey page compact but unchanged; sustained analysis uses
% a slightly smaller reading size so longer patristic dossiers remain usable.
\newenvironment{deepstudy}
  {\begingroup\small}
  {\par\endgroup}

\newenvironment{threestable}[3]
{\small\begin{longtable}{>{\bfseries\raggedright\arraybackslash}p{0.18\linewidth}>{\raggedright\arraybackslash}p{0.34\linewidth}>{\raggedright\arraybackslash}p{0.40\linewidth}}
\toprule
\textbf{#1} & \textbf{#2} & \textbf{#3}\\
\midrule
\endhead}
{\bottomrule\end{longtable}}

\newenvironment{fourstable}[4]
{\footnotesize\begin{longtable}{>{\bfseries\raggedright\arraybackslash}p{0.15\linewidth}>{\raggedright\arraybackslash}p{0.23\linewidth}>{\raggedright\arraybackslash}p{0.25\linewidth}>{\raggedright\arraybackslash}p{0.27\linewidth}}
\toprule
\textbf{#1} & \textbf{#2} & \textbf{#3} & \textbf{#4}\\
\midrule
\endhead}
{\bottomrule\end{longtable}}

\newtcolorbox{studybox}[1]{
  enhanced,
  breakable,
  colback=white,
  colbacktitle=white,
  coltitle=black,
  frame hidden,
  boxrule=0pt,
  sharp corners,
  borderline west={1.2pt}{0pt}{black},
  left=8pt,
  right=2pt,
  top=2pt,
  bottom=2pt,
  toptitle=0pt,
  bottomtitle=2pt,
  before skip=8pt,
  after skip=8pt,
  title={#1},
  fonttitle=\small\bfseries
}

\newtcolorbox{massbox}[1]{
  enhanced,
  breakable,
  colback=white,
  colframe=black,
  colbacktitle=white,
  coltitle=black,
  boxrule=0.55pt,
  sharp corners,
  left=7pt,
  right=7pt,
  top=4pt,
  bottom=4pt,
  toptitle=3pt,
  bottomtitle=3pt,
  titlerule=0.45pt,
  before skip=8pt,
  after skip=8pt,
  title={#1},
  fonttitle=\small\bfseries
}

\newtcolorbox{sourcecard}[1]{
  enhanced,
  colback=white,
  colframe=black,
  colbacktitle=white,
  coltitle=black,
  boxrule=0.4pt,
  sharp corners,
  left=5pt,
  right=5pt,
  top=4pt,
  bottom=4pt,
  toptitle=3pt,
  bottomtitle=3pt,
  titlerule=0.4pt,
  before skip=4pt,
  after skip=4pt,
  title={#1},
  fonttitle=\small\bfseries
}

\tikzset{
  guide node/.style={draw=black, fill=white, align=center, inner sep=5pt, font=\small},
  guide emphasis/.style={draw=black, very thick, fill=white, align=center, inner sep=5pt, font=\small},
  guide arrow/.style={-{Stealth[length=2mm]}, thick, draw=black}
}

\input{liturgy/roman-rite/1962/reference/altar-server-guides/shared/server-guide-format}
\SolemnMassGuidetrue
\addtolength{\textheight}{1pt}

\newcommand{\GuideRoot}{liturgy/roman-rite/1962/reference/altar-server-guides/03-solemn-mass}
\hypersetup{
  pdftitle={Serving Solemn Mass},
  pdfsubject={A complete model-roster action, response, memorization, and Roman-pronunciation guide for young servers of the 1962 Roman Missal}
}

\begin{document}
\input{liturgy/roman-rite/1962/reference/altar-server-guides/shared/sung-mass-diagrams}
\input{liturgy/roman-rite/1962/reference/altar-server-guides/03-solemn-mass/sections/00-start-here}
\serverguidetitle
  {Serving Solemn Mass}
  {Know your part. Move with the team.}
  {Responses, pronunciation, roles, and actions for servers ages 8--14}

\begin{deepstudy}
\input{liturgy/roman-rite/1962/reference/altar-server-guides/03-solemn-mass/sections/00-begin}
\input{liturgy/roman-rite/1962/reference/altar-server-guides/shared/solemn-role-sheets}

\appendix
\input{liturgy/roman-rite/1962/reference/altar-server-guides/shared/solemn-ceremony}
\input{liturgy/roman-rite/1962/reference/altar-server-guides/shared/response-reference}

\newcommand{\GuideFormStatement}{Non-pontifical Solemn Mass with celebrant,
deacon, and subdeacon, according to the Order and rubrics of the 1962 Roman
Missal.}
\newcommand{\GuideBoundaryStatement}{The guide accounts for the complete
ordinary route of its declared lay-server roster. Sacred-minister and choir
acts appear only as cues or handoff points; their offices are not reassigned to
young servers. Special rites require their own sheets.}
\newcommand{\GuideRosterStatement}{The core model uses an MC, thurifer, boat
bearer, two acolytes, and four torchbearers. A processional cross-bearer is an
optional local addition, not part of the core roster. This teaching route
requires the separate thurifer and boat-bearer roles; it states how the
optional cross-bearer rejoins the route. Local traffic remains subject to the
priest and MC.}
\newcommand{\GuideSafetyStatement}{A young server needs supervised,
site-specific practice before handling flames, hot charcoal, vessels, a
Communion plate, or unfamiliar steps.}
\newcommand{\GuideCeremonialReferences}{A.\ Fortescue and J.\ B.\ O'Connell,
\textit{The Ceremonies of the Roman Rite Described}, 12th revised edition
(1962), names the selected ceremonial family. The directly read 1917 edition
and the Archconfraternity of Saint Stephen's 2022 handbook were used only as
subordinate traffic and discrepancy controls; the 1962 Missal overrides both.}
% The learner route already entered appendix mode before the planning and
% trainer material. Do not let the shared terminal apparatus restart its
% appendix lettering.
\let\appendix\relax
% Keep the terminal heading with the opening scope paragraph.
\Needspace{0.16\textheight}
\input{liturgy/roman-rite/1962/reference/altar-server-guides/shared/terminal-apparatus}
{\scriptsize\input{liturgy/roman-rite/1962/reference/altar-server-guides/03-solemn-mass/generation-metadata}}
\end{deepstudy}
\end{document}
